\section{Atomwise invertible multiplier linearization}
\label{sec:linearization}

\subsection{Integer provenance at the branch level}

\begin{lemma}[Canonical integer factorization]
\label{lem:correct-integer-factorization}
Write
\(\sigma_\beta=g_\beta\sigma_{\beta,0}\).  Then
\begin{equation}
 \frac{\Omega_\beta}{c_\beta}
 =
 \kappa_\beta\ell_\beta v_\beta\sigma_{\beta,0}
 =
 \kappa_\beta\ell_\beta v_\beta
 \frac{\sigma_\beta}{g_\beta}
 \in\mathbb Z.
\label{eq:correct-integer-factorization}
\end{equation}
\end{lemma}

\begin{proof}
Since
\(b_\beta=g_\beta B_\beta=c_\beta\kappa_\beta\),
\[
 \Omega_\beta
 =
 \ell_\beta v_\beta
 (g_\beta\sigma_{\beta,0})B_\beta
 =
 c_\beta\kappa_\beta\ell_\beta
 v_\beta\sigma_{\beta,0}.
\]
\end{proof}

\begin{example}[No integer quotient \(\ell vB/c\)]
\label{ex:no-integer-coefficient}
Take
\[
 c=2,\quad\kappa=1,\quad b=2,\quad
 \sigma=2,\quad g=2,\quad B=1,\quad\ell=v=1.
\]
Then \(\Omega/c=1\), but \(c\nmid\ell vB\).  Therefore
\((\ell vB/c)\sigma\) has the correct rational value but does not
exhibit an integer coefficient.  Division by \(c\) below is a
finite-field inversion, not this extra integer divisibility.
\end{example}

\subsection{The actual opened value is the invertible divisor}

\begin{theorem}[Atomwise affine multiplier formula]
\label{thm:atomwise-linearization}
Let \(p=(\beta,z)\) be a nonconstant opened atom with
\(q_X\nmid u_p\).  Define
\begin{equation}
 A_p
 =
 \varepsilon_\beta\ell_\beta v_\beta B_\beta
 (c_\beta u_p)^{-1}
 \in\mathbb F_{q_X}^{\times}.
\label{eq:atom-A}
\end{equation}
Then
\begin{equation}
 \boxed{
 \omega_\beta
 =
 A_p(m_pj_p+h_0)
 \quad\text{in }\mathbb F_{q_X}^{\times}.}
\label{eq:atom-affine-omega}
\end{equation}
Equivalently,
\begin{equation}
 A_p
 =
 \varepsilon_\beta\kappa_\beta\ell_\beta v_\beta
 (u_pg_\beta)^{-1}.
\label{eq:atom-A-second}
\end{equation}
\end{theorem}

\begin{proof}
All factors in \eqref{eq:atom-A} are nonzero modulo \(q_X\).
Reduce the branch multiplier and use the atom source identity:
\begin{align*}
 \omega_\beta
 &\equiv
 \varepsilon_\beta\ell_\beta v_\beta B_\beta
 c_\beta^{-1}\sigma_\beta\\
 &\equiv
 \varepsilon_\beta\ell_\beta v_\beta B_\beta
 (c_\beta u_p)^{-1}(m_pj_p+h_0).
\end{align*}
This proves \eqref{eq:atom-affine-omega}.  The identity
\(B_\beta g_\beta=c_\beta\kappa_\beta\) gives
\eqref{eq:atom-A-second}.
\end{proof}

\begin{remark}[Dependence on \(z\)]
The left side \(\omega_\beta\) is constant for
\(z\in I_\beta^\ast\).  The coefficient \(A_p\) and source row
\(m_p\) generally vary with \(z\), because
\[
 u_p=U_\theta(\tau_\beta+B_\beta z),\qquad
 m_p=M_\theta(\tau_\beta+B_\beta z).
\]
No theorem here treats \(u\) or \(m\) as fixed on \(\beta\).
\end{remark}

\begin{theorem}[Fine-cell injectivity]
\label{thm:fine-cell-injectivity}
Let \(C\) be a fine provenance cell consisting of atoms with
\(q_X\nmid u_p\).  Then the common values in
\cref{thm:lossless-atom-cells} define
\[
 A_C=
 \varepsilon_C\ell_Cv_CB_C(c_Cu_C)^{-1},
\]
and
\begin{equation}
 p\longmapsto\omega_{\beta(p)}
 =
 A_C(m_Cj_p+h_0)
\label{eq:fine-cell-map}
\end{equation}
is injective on \(P_C\).
\end{theorem}

\begin{proof}
By \cref{thm:lossless-atom-cells}(ii), \(P_C\) is already a
singleton.  Equivalently, one may verify the affine implication
directly as follows.
If two atoms \(p,p'\in P_C\) have equal multipliers, then
\[
 A_Cm_C(j_p-j_{p'})=0\pmod{q_X}.
\]
Since \(A_Cm_C\ne0\), the two orbit coordinates are congruent
modulo \(q_X\).  They lie in the true orbit window of diameter below
\(q_X\), so \(j_p=j_{p'}\).  By
\cref{thm:lossless-atom-cells}(i), the common cell record and this
coordinate reconstruct the same \((\beta,z)\).  Hence \(p=p'\).
\end{proof}

\begin{remark}
This theorem records compatibility of the affine formula with the
lossless provenance partition.  It supplies no nontrivial
within-cell dispersion because the maximally literal cells are
singletons.
\end{remark}

\begin{remark}[Same branch, different atoms]
If \(z\ne z'\) belong to one branch \(\beta\), then both atoms have
the same \(\omega_\beta\), but their opened source values generally
differ.  They therefore lie in different fine cells, and their
collision is retained as a cross-cell term.  Fine-cell injectivity
does not convert a branchwise point mass into dispersed mass.
\end{remark}
